\section{Basic Exercises}

Determine the symmetric matrix $A$ such that $Q(\vx) = \vx^T A \vx$.
\begin{tasks}[style=enumerate](1)
\task
  \begin{displaymath}
    Q([x_1 \ x_2 \ x_3]^T)
    = 3x_1^2 + 4 x_1 x_2 + 5 x_2 x_1 - x_1 x_3
  \end{displaymath}
\task
  \begin{displaymath}
    Q([x_1 \ x_2 \ x_3]^T)
    = 3x_1^2 + 3x_2^2 + 5x_3^2 + 6x_1 x_2 + 2x_1 x_3 + 2x_2 x_3
  \end{displaymath}
\end{tasks}
Determine the quadratic form $Q(\vx)$ such that $Q(\vx) = \vx^T A
\vx$.
\begin{tasks}[resume, style=enumerate](1)
\task
  \begin{align*}
    A =
    \begin{bmatrix}
      2 & 0 & 0 \\
      0 & 9 & 3 \\
      0 & 3 & 1
    \end{bmatrix}
  \end{align*}
\end{tasks}
Determine the type of the quadratic form $Q(\vx) = \vx^T A \vx$.
\begin{tasks}[resume, style=enumerate](1)
\task
  \begin{displaymath}
    A =
    \begin{bmatrix}
      2 & 0 & 0 \\
      0 & 9 & 3 \\
      0 & 3 & 1
    \end{bmatrix}
  \end{displaymath}
\end{tasks}
Determine the singular values of the following matrices.
\begin{tasks}[resume, style=enumerate](1)
\task
  \begin{displaymath}
    \begin{bmatrix}
      1 & 1 & 0 & 2 \\
      3 & 1 & 1 & -2
    \end{bmatrix}
  \end{displaymath}
\end{tasks}
Determine a singular value decomposition for each of the following
matrices.
\begin{tasks}[resume, style=enumerate](1)
\task
  \begin{displaymath}
    \begin{bmatrix}
      3 & 2 \\
      2 & 3 \\
      2 & -2
    \end{bmatrix}
  \end{displaymath}
\end{tasks}

\section{True/False}

\trueFalseHeader

\begin{tasks}[resume, style=enumerate](1)
\task
  Every positive-definite symmetric matrix is invertible.
\task
  For any matrix $A$ and quadratic form $Q(\vx)$, if $Q(\vx) =
  \vx^TA\vx$, then $A$ is symmetric.
\task
  For any quadratic form $Q(\vx)$, the vector $\argmax_{\|\vx\| = 1}
  Q(\vx)$ is unique.
\task
  $\|\vx\|^2$ is a quadratic form.
\task
  A positive definite quadratic form $Q(\vx)$ satisfies $Q(\vx) > 0$
  for all $\vx \in \R^n$.
\task
  $\vx^T A \vx$ defines a quadratic form only if $A$ is symmetric.
\task
  An orthogonal diagonalization of a symmetric matrix $A = PDP^T$ is
  also a singular value decomposition of $A$.
\task
  If $A$ is positive semidefinite, then every eigenvalue of $A$ is
  also a singular value of $A$.
\end{tasks}

\section{More Difficult Problems}

\begin{tasks}[resume, style=enumerate](1)
\task
  Consider the quadratic form $Q(\vx) = \vx^T A \vx$ where $A$ is
  defined as follows.
  \begin{displaymath}
    \begin{bmatrix}
      1 & 1 & 5 \\
      1 & 5 & 1 \\
      5 & 1 & 1
    \end{bmatrix}
  \end{displaymath}
  Determine the following.
  \begin{enumerate}[label=(\alph*)]
  \item
    $\min_{\|\vx\| = 1} Q(\vx)$
  \item
    $\max_{\|\vx\| = 1}  Q(\vx)$
  \item
    $\argmin_{\|\vx\| = 1} Q(\vx)$
  \item
    $\argmax_{\|\vx\| = 1} Q(\vx)$
  \end{enumerate}
  \textit{Hint:} $4$ is an eigenvalue of $A$.
\task
  \faCalculator \ Repeat the previous problem with the following
  quadratic form.
  \begin{displaymath}
    Q([x_1 \ x_2 \ x_3]^T) =
    3x_1^2 + 3x_2^2 + 5x_3^2 + 6x_1 x_2 + 2x_1 x_3 + 2x_2 x_3
  \end{displaymath}
\task
Consider the following matrices and vector.
  \begin{displaymath}
    A =
    \begin{bmatrix}
      2 & -4 & 4 \\
      2 & 2 & 1 \\
      -2 & 4 & -4 \\
    \end{bmatrix}
    \qquad
    AA^T =
    \begin{bmatrix}
      36 & 0 & -36 \\
      0 & 9 & 0 \\
      -36 & 0 & 36
    \end{bmatrix}
    \qquad
    \vb =
    \begin{bmatrix}
      10 \\ 18 \\ -14
    \end{bmatrix}
  \end{displaymath}
  \begin{enumerate}[label=(\alph*)]
  \item
    Determine a \textit{reduced} singular value decomposition of $A$,
    given that $6\sqrt{2}$ and $3$ are its nonzero singular values and
    $\rank(A) = 2$.
  \item
    Determine the length of the \textit{shortest} least-squares
    solution to the equation $A\vx = \vb$.  \textit{Hint:} Use the
    pseudoinverse of $A$.
  \end{enumerate}
\task
  \label{prob:svd-prev}
  Consider the following matrices.
  \begin{align*}
    A =
    \begin{bmatrix}
      2 & 2 & 2 \\
      1 & 1 & -2 \\
      2 & 2 & 2
    \end{bmatrix}
    \qquad
    V =
    \begin{bmatrix}
      \sqrt{3} / 3 & \sqrt{6} / {6} \\
      \sqrt{3} / 3 & \sqrt{6} / {6} \\
      \sqrt{3} / 3 & -\sqrt{6} / {3}
    \end{bmatrix}
  \end{align*}
  \begin{enumerate}[label=(\alph*)]
  \item
    Without first determining $A^TA$ or $AA^T$, determine left singular vectors
    of $A$ associated with nonzero singular values, i.e., determine the
    columns of $U$ in the \textit{reduced} SVD of $A$, where $U\Sigma
    V^T$ is the reduced SVD of $A$.
  \item
    Without first determining $A^TA$ or $AA^T$, determine the singular
    values of $A$.
  \end{enumerate}
\task
  Consider the following reduced singular value decomposition of a
  matrix $A$.
  \begin{align*}
    A
    =
    \begin{bmatrix}
      \sqrt{3} / 3 & -\sqrt{6} / 6 \\
      \sqrt{3} / 3 & \sqrt{6} / 3 \\
      \sqrt{3} / 3 & -\sqrt{6} / 6
    \end{bmatrix}
    \begin{bmatrix}
      3 & 0 \\
      0 & 2
    \end{bmatrix}
    \begin{bmatrix}
      \sqrt{14} / 7 & \sqrt{14} / 14 & 3 \sqrt{14} / 14 \\
      0 & -3 \sqrt{10} / 10 & \sqrt{10} / 10
    \end{bmatrix}
  \end{align*}
  \begin{enumerate}[label=(\alph*)]
  \item
    Determine a basis for $\col A$.
  \item
    Determine a basis for $\nul A$.
  \item
    Determine an orthogonal diagonalization of $A^T A$.
  \end{enumerate}
\end{tasks}

\section{Challenge Problems}

\begin{tasks}[resume, style=enumerate](1)
\task
  Let $A$ be as in Problem \ref{prob:svd-prev} above.  Determine the matrix that
  implements orthogonal projection onto $\col A$.
\end{tasks}
